% © 2026 Ing. David Jaroš — CC BY-NC-ND 4.0
%
% File: research_tracks/EW/weinberg_angle_mz_bridge_closure.tex
% Status: RESEARCH TRACK / BRIDGE-PROVEN.
% Purpose: Combine UBT Weinberg boundary, odd-spinor beta coefficients, scale and alpha/QED bridge.

% UBT-AI-PROVENANCE-BEGIN
% schema: ubt-ai-provenance/v1
% tier: C_working
% ai_assistance: disclosed
% human_review: risk-based
% editorial_responsibility: Ing. David Jaroš
% policy: ../../AI_PROVENANCE.md
% notice: Working material; exhaustive human review is not claimed.
% UBT-AI-PROVENANCE-END

\documentclass[12pt,a4paper]{article}
\usepackage[a4paper,margin=2.5cm]{geometry}
\usepackage{amsmath,amssymb,mathtools,booktabs}
\usepackage[most]{tcolorbox}
\usepackage{hyperref}
\usepackage{xcolor}
\usepackage{microtype}

\hypersetup{colorlinks=true,linkcolor=blue!60!black,urlcolor=blue!60!black,citecolor=green!50!black}

\newtcolorbox{statusbox}[1][]{
  colback=green!6!white,
  colframe=green!45!black,
  arc=3pt,
  boxrule=1pt,
  title=Bridge-Proof Status,
  #1
}

\title{Z-Pole Weinberg Angle Closure on the UBT Odd-Spinor Branch}
\author{Ing. David Jaroš}
\date{2026-06-14}

\begin{document}
\maketitle

\begin{abstract}
This note combines the UBT hypercharge boundary
$\sin^2\theta_W(M_{\rm UBT})=3/8$, the odd-spinor threshold beta coefficients
$(b_1,b_2)=(33/5,1)$, the compact unification scale
$M_{\rm UBT}=2\times10^{16}$ GeV, and the alpha/QED bridge value
$\alpha_{\rm EM}^{-1}(M_Z)=127.934499434\ldots$.  The resulting one-loop
closure gives
\[
  \sin^2\theta_W(M_Z)=0.23122\ldots.
\]
The status is bridge-proven: the numerical descent is proven relative to the
current UBT bridges for the odd-spinor threshold spectrum, compact scale, and
QED running of the electromagnetic coupling.
\end{abstract}

\begin{statusbox}
\textbf{Claim.}  Within the current UBT bridge assumptions,
\[
  \boxed{\sin^2\theta_W(M_Z)=0.23122\ldots}
\]
follows from the $3/8$ boundary.  This is an internal UBT bridge-proof status,
not a claim of external peer-reviewed acceptance.
\end{statusbox}

\section{Inputs from previous UBT results}

The proof uses:
\begin{align}
  \sin^2\theta_W(M_{\rm UBT})&=\frac38,\\
  b_1&=\frac{33}{5},\\
  b_2&=1,\\
  M_{\rm UBT}&=2\times10^{16}\,\mathrm{GeV},\\
  M_Z&=91.1876\,\mathrm{GeV},\\
  \alpha_{\rm EM}^{-1}(M_Z)&=127.934499434\ldots.
\end{align}
The first input is the UBT hypercharge norm theorem.  The beta coefficients are
proved in \texttt{weinberg\_angle\_odd\_spinor\_threshold\_beta\_theorem.tex}.
The electromagnetic coupling is the $Z$-pole value obtained by the alpha/QED
running bridge.  The compact scale is the UBT/GUT-scale bridge used in the
EW1+RG branch.

\section{Algebraic closure}

Let
\[
  L=\frac{1}{2\pi}\ln\frac{M_{\rm UBT}}{M_Z}.
\]
For the numerical values above,
\[
  L=5.255549236755\ldots.
\]
The GUT-normalised one-loop running is
\[
  \alpha_i^{-1}(M_Z)=A+b_iL,
\]
where
\[
  A:=\alpha_{\rm UBT}^{-1}.
\]
The physical hypercharge coupling satisfies
\[
  \alpha_Y^{-1}=\frac53\alpha_1^{-1}.
\]
Therefore
\[
  \alpha_{\rm EM}^{-1}(M_Z)
  =
  \alpha_Y^{-1}(M_Z)+\alpha_2^{-1}(M_Z)
  =
  \frac53(A+b_1L)+(A+b_2L).
\]
Using $b_1=33/5$, $b_2=1$, this becomes
\[
  \alpha_{\rm EM}^{-1}(M_Z)
  =
  \frac83A+12L.
\]
Thus
\[
  A
  =
  \frac38\left[\alpha_{\rm EM}^{-1}(M_Z)-12L\right].
\]
Substitution gives
\[
  A=24.3254657224\ldots.
\]
Then
\[
  \alpha_2^{-1}(M_Z)=A+L=29.5810149592\ldots.
\]
Since
\[
  \sin^2\theta_W(M_Z)=\frac{\alpha_{\rm EM}(M_Z)}{\alpha_2(M_Z)}
  =
  \frac{\alpha_2^{-1}(M_Z)}{\alpha_{\rm EM}^{-1}(M_Z)},
\]
we obtain
\[
  \sin^2\theta_W(M_Z)
  =
  \frac{29.5810149592\ldots}{127.934499434\ldots}
  =
  0.23122\ldots.
\]

\section{Interpretation}

The descent can be summarised as
\[
  \frac38
  \xrightarrow[\alpha_{\rm EM}^{-1}(M_Z),\ M_{\rm UBT}]
  {b_1=33/5,\ b_2=1}
  0.23122\ldots.
\]
Unlike the previous provisional RG branch, the beta coefficients are now
computed from the odd-spinor threshold spectrum rather than inserted as a loose
MSSM analogy.

\section{Remaining external-facing caveat}

This is a bridge proof, not an experimental publication claim.  A reviewer may
still challenge the bridge assumptions:
\begin{enumerate}
  \item the derivation of the compact scale $M_{\rm UBT}$;
  \item the alpha/QED running bridge to $\alpha_{\rm EM}^{-1}(M_Z)$;
  \item the identification of the odd-spinor threshold spectrum with the active
  one-loop states between $M_Z$ and $M_{\rm UBT}$.
\end{enumerate}
Within those stated UBT bridges, however, the Weinberg angle is closed at the
same internal proof level as the alpha chain.

\end{document}
